To separate model-dependent force response from trajectory-level stochasticity, we evaluated the same 639 failure-attribution structures with twelve MACE checkpoints spanning three from-scratch models, two fine-tuned families at three training seeds each, and three foundation-model representatives. The structures were drawn from nine trajectory/seed branches and intentionally stratified across early-stable, pre-onset, precursor and threshold regimes. The resulting set therefore represents the audited failure-attribution frame rather than a time-weighted molecular-dynamics ensemble. All 7,668 paired evaluations were finite and complete.

The models agreed closely on early-stable and stable pre-onset structures, with median maximum forces of approximately 2.4--3.2~eV~\AA$^{-1}$. Model divergence increased in the precursor regime: the median maximum force ranged from 3.2~eV~\AA$^{-1}$ for the baseline checkpoint to 29.9--30.8~eV~\AA$^{-1}$ for the two foundation medium models. At the threshold, the median twelve-model force range reached 3,717.51~eV~\AA$^{-1}$, compared with 3.07~eV~\AA$^{-1}$ for early-stable controls. These values identify a reproducible model-divergence signature around the apparent instability, while remaining agnostic about which force is physically correct.

The candidate structures also occupy a shifted region of the stored SOAP-PCA diagnostic. Every candidate stage lies outside the training q99 envelope; the candidate PCA coordinates are separated from both the training/validation/test region and the NEB reference region. This establishes that the candidates are deliberately selected out-of-distribution probes relative to the current reference embedding, not unbiased samples of the original DFT manifold. Because no reference DFT forces are available for the 2,050-atom cells, these static comparisons do not by themselves establish force error or causality. Targeted DFT calculations on reduced, PBC-preserving fragments are therefore required for physical validation.
